\section{The scalar density of the original phase graph source}
\label{sec:pgd-density}

Retain the actual quartet, its full order, its full arithmetic unit and its
original primitive from MBG.1--MBG.3. In particular
\[
 \begin{gathered}
 h(s)=\prod_{r\in\{\rho,\bar\rho,1-\rho,1-\bar\rho\}}(s-r)^m,
 \quad \rho=\tfrac12+\delta+i\gamma,\quad
 0<\delta<\tfrac12,\quad \gamma>2,\quad m\in\mathbb Z_{>0},\\
 g=2\xi,\quad v_h=g/h,\quad \Phi'=h,\quad\Phi(0)=0,\quad
 \Phi(\tfrac12+it)=A+iq_\Phi(t),\quad A>0,\\
 q_\Phi'(t)=
 \left[((t-\gamma)^2+\delta^2)((t+\gamma)^2+\delta^2)\right]^m,
 \quad q_\Phi(0)=0.
 \end{gathered}
 \tag{PGD.1}\label{eq:pgd-data}
\]
Thus the actual polynomial \(q_\Phi\) is real and odd of degree
\(p_\Phi=4m+1\), with leading coefficient \(1/(4m+1)\).
These facts follow by differentiating the retained primitive on its
original line, not by choosing a different phase. Put
\[
 \begin{gathered}
 w(t)=\frac{|v_h(1/2+it)|^2}{2\pi},\quad
 m_k=w^{*k},\quad \mu_h=\int_{\mathbb{R}}w(t)\,dt>0,\\
 s_j=\tfrac12+it_j,\quad S=\tfrac k2+iu,\quad u=\sum_jt_j,\\
 (\mathcal V_NP)(t_1,\ldots,t_k)
       =\prod_jv_h(s_j)P\left(\sum_js_j\right),\qquad
 \mathsf T_k=\sum_j\Phi(s_j)=kA+i\sum_jq_\Phi(t_j),\\
 M_N=\mathcal V_N^*\mathcal V_N,\qquad
 W_N=\mathcal V_N^*(I+\mathsf T_k^*\mathsf T_k)\mathcal V_N,
 \quad \mathcal P_N=\mathbb{C}[S]_{\leq N}.
 \end{gathered}
 \tag{PGD.2}\label{eq:pgd-source}
\]
The Hilbert measure before multiplication by the displayed full units is
\(\prod_jdt_j/(2\pi)\). Consequently all following scalar convolutions
use the already displayed density \(w(t)dt\); they introduce no further
factor of \(2\pi\). The adjoints use the original increasing-power
coefficient frame. These are the full phase graph forms; the subtracted
finite-jet boundary in MBG.24 is a different, explicitly related map.

\subsection{Exact fibre and three-convolution formulas}

For \(k\geq2\), eliminate the last variable by
\(t_k=u-\sum_{j<k}t_j\), and define
\[
 r_k(u)=\int_{\mathbb{R}^{k-1}}
 \left(1+k^2A^2+
           \left[\sum_{j=1}^kq_\Phi(t_j)\right]^2\right)
 \prod_{j=1}^kw(t_j)\,dt_1\cdots dt_{k-1}.
 \tag{PGD.3}\label{eq:pgd-fibre}
\]
The linear substitution from \((t_1,\ldots,t_k)\) to
\((t_1,\ldots,t_{k-1},u)\) has determinant one. Tonelli's theorem
on the nonnegative diagonal integrand, followed by polarization for
cross-pairings, therefore proves
\[
 \langle P,W_NQ\rangle
 =\int_{\mathbb{R}}\overline{P(k/2+iu)}Q(k/2+iu)r_k(u)\,du.
 \tag{PGD.4}\label{eq:pgd-scalar-form}
\]
Here \(\langle P,W_NQ\rangle\) denotes the original coefficient pairing.
All polynomial moments converge absolutely by the retained exponential
decay of \(w\), proved quantitatively below. This also justifies each
finite expansion used next. In the original density notation the exact
formula is
\[
 \boxed{\quad
 r_k=(1+k^2A^2)w^{*k}
       +k(q_\Phi^2w)*w^{*(k-1)}
       +k(k-1)(q_\Phi w)*(q_\Phi w)*w^{*(k-2)}.
 \quad}
 \tag{PGD.5}\label{eq:pgd-three-convolutions}
\]
For this displayed formula \(k\geq2\), and \(w^{*0}\) is the unit
Dirac measure \(\delta_0\) for convolution, with mass one. For \(k=1\)
there is no third term and the complete formula is explicitly
\[
 r_1(u)=[1+A^2+q_\Phi(u)^2]w(u).
 \tag{PGD.6}\label{eq:pgd-one-factor}
\]
Indeed the square in (PGD.3) has exactly \(k\) diagonal terms and
\(k(k-1)\) ordered terms with unequal indices. Fubini and the same
determinant-one substitution give their three convolutions. The third
convolution in (PGD.5) is a signed function: \(q_\Phi w\) is odd,
and it must not be replaced by its absolute value in that equality.

Conjugation and the full quartet symmetry give \(w(-t)=w(t)\),
whereas \(q_\Phi\) is odd. Hence
\(\int q_\Phi w=0\). Integrating all three terms retains the exact
source mass:
\[
 \int r_k(u)\,du
  =(1+k^2A^2)\mu_h^k
        +k\mu_h^{k-1}\int q_\Phi(t)^2w(t)\,dt.
 \tag{PGD.7}\label{eq:pgd-mass}
\]
All functions \(r_k\) are even. The fibre representation proves
\[
 (1+k^2A^2)m_k(u)\leq r_k(u)
 \leq (1+k^2A^2)m_k(u)+k^2z_k(u),\qquad
 z_k=(q_\Phi^2w)*w^{*(k-1)}.
 \tag{PGD.8}\label{eq:pgd-density-order}
\]
The upper inequality follows from
\((\sum_jx_j)^2\leq k\sum_jx_j^2\) for real \(x_j\), followed
by permutation of the \(k\) original integration variables.

The equality cases at fixed ordinate can also be calculated. For two
factors and \(u=0\), oddness gives
\(q_\Phi(t)+q_\Phi(-t)=0\) at every point of the fibre, and hence
\[
 r_2(0)=(1+4A^2)m_2(0),\qquad
 [(q_\Phi w)*(q_\Phi w)](0)
             =-[(q_\Phi^2w)*w](0).
 \tag{PGD.9}\label{eq:pgd-two-zero}
\]
For \(k=2\) and \(u\ne0\), the polynomial
\(q_\Phi(x)+q_\Phi(u-x)\) has leading term
\(u x^{p_\Phi-1}\), so is nonconstant. For \(k\geq3\) and
every \(u\), the fibre polynomial is nonconstant: set
\(t_1=x,t_2=y,t_3=u-x-y,t_j=0\) for \(j>3\).
In its homogeneous part of degree \(p_\Phi\), the coefficient
of \(x^{p_\Phi-1}y\) is \(-1\). No lower phase coefficient
can cancel it. The product density on the fibre is positive almost
everywhere, since each factor is the squared modulus of a nonzero
entire function evaluated on a real affine coordinate. A nonzero
polynomial cannot vanish on an open set; moreover its zero set has
Lebesgue measure zero, as follows inductively by viewing it as a
one-variable polynomial with polynomial coefficients and applying
Fubini. Its squared integral against this density is therefore
strictly positive. This proves
\[
 r_k(u)>(1+k^2A^2)m_k(u)\quad
 \begin{cases}
  k=2,\ u\ne0,\\
  k\geq3,\ u\in\mathbb{R}.
 \end{cases}
 \tag{PGD.10}\label{eq:pgd-strict-fibre}
\]
For one factor, (PGD.6) states the exact equality cases, including
zeros of the actual \(w\). The preceding statements are about the
whole phase-square fibre, not a sign assertion for its individual
mixed convolution in (PGD.5).

\subsection{Uniform actual tail constants and the phase insertion}

The original BT12--BT16 proof supplies the following constants and
bound. They are retained explicitly to specify the estimate being used:
\[
 \begin{gathered}
 \alpha=\pi/2,\quad p_0=8m-9/2>1,\quad
 R_0=(\delta^2+\gamma^2)^{1/2},\quad T_h=\max(1,2R_0),\\
 Z_*=2+\sqrt{5/2},\quad
 c_\Gamma=\sqrt2e^{-7/6},\quad
 C_\Gamma=\sqrt{2\sqrt5}\,e^{2/3},\quad \beta_0=8m-5,\\
 C_{\rm comp}=\frac{C_\Gamma}{\sqrt\pi}(T_h^2+1/4)^2
       [1+2\sqrt{T_h^2+1/4}]^2
       \delta^{-8m}(1+T_h)^{\beta_0},\\
 C_{\rm tail}=\frac{25C_\Gamma Z_*^2}{16\sqrt\pi}
       (4/3)^{4m}(1+1/T_h)^{\beta_0},\qquad
 C_h=\max(C_{\rm comp},C_{\rm tail}),\\
 w(t)\leq C_he^{-\alpha|t|}(1+|t|)^{-p_0},\qquad
 M_\alpha=\int e^{\alpha t}w(t)\,dt
         =\int e^{-\alpha t}w(t)\,dt\in(0,\infty).
 \end{gathered}
 \tag{PGD.11}\label{eq:pgd-original-tail}
\]
These are the original arithmetic bounds proved from the floor formula
for zeta and the complete Binet identity in BT. In particular their
use does not assume a positive lower bound at individual zeta zeros.
They imply \(M_\alpha\leq2C_h/(p_0-1)\).
Write the complete phase as
\(q_\Phi(t)=\sum_{j=0}^{p_\Phi}q_jt^j\), and set
\[
 \begin{gathered}
 Q_h=\sum_{j=0}^{p_\Phi}|q_j|>0,\qquad
 a=13/2,\qquad C_\Phi=C_hQ_h^2H_a,\\
 H_a=\sup_{x\geq0}(1+x)^ae^{-2\alpha x}
  =\begin{cases}
      1,&a\leq2\alpha,\\
      e^{2\alpha-a}(a/(2\alpha))^a,&a>2\alpha.
    \end{cases}
 \end{gathered}
 \tag{PGD.12}\label{eq:pgd-phase-constants}
\]
Every \(q_j\) here is the coefficient of the original polynomial in
(PGD.1). The phase has not been replaced by the polynomial with
absolute coefficients; that polynomial is used only for the bound
\(|q_\Phi(t)|\leq Q_h(1+|t|)^{p_\Phi}\).
Because \(2p_\Phi-p_0=13/2\), (PGD.11) gives
\[
 q_\Phi(t)^2w(t)
       \leq C_hQ_h^2e^{-\alpha|t|}(1+|t|)^a.
 \tag{PGD.13}\label{eq:pgd-inserted-tail}
\]
The formula for \(H_a\) follows by differentiating
\(a\log(1+x)-2\alpha x\), whose derivative is strictly decreasing.
Let \(t_-:=\max(-t,0)\), and retain the actual finite one-sided moment
\[
 I_a^-:=\int (1+t_-)^ae^{\alpha t}w(t)\,dt>0.
 \tag{PGD.14}\label{eq:pgd-negative-moment}
\]
It is finite with the explicit bound
\[
 I_a^-\leq C_h\left[
       \frac1{p_0-1}
       +\sum_{j=0}^{7}\binom7j\frac{j!}{(2\alpha)^{j+1}}
                           \right].
 \tag{PGD.15}\label{eq:pgd-negative-moment-bound}
\]
Indeed on the positive half-line the integrand is at most
\(C_h(1+t)^{-p_0}\). On the negative half-line put \(x=-t\);
its bound is \(C_he^{-2\alpha x}(1+x)^{a-p_0}\), which is at
most \(C_he^{-2\alpha x}(1+x)^7\). Expanding this last polynomial
and integrating each monomial gives (PGD.15). No positive-part
moment of order \(a\) at the exponential boundary is required.

For every \(k\geq2\), the actual inserted convolution obeys
\[
 \boxed{\quad
 z_k(u)\leq C_\Phi(k-1)^aM_\alpha^{k-2}I_a^-
               e^{-\alpha|u|}(1+|u|)^a.
 \quad}
 \tag{PGD.16}\label{eq:pgd-z-envelope}
\]
Here is its proof with the one-sided tilt preserved. Put
\(b(t)=e^{\alpha t}w(t)\), of mass \(M_\alpha\), and
\(f(t)=e^{\alpha t}q_\Phi(t)^2w(t)\).
For \(t\geq0\), (PGD.13) gives
\(f(t)\leq C_\Phi(1+t)^a\). For \(t<0\), it gives
\(f(t)\leq C_hQ_h^2H_a\). Thus on both half-lines
\(f(t)\leq C_\Phi(1+t_+)^a\).
For \(u\geq0\) and every \(v\),
\(1+(u-v)_+\leq(1+u)(1+v_-)\). Consequently, with \(n=k-1\),
\[
 \begin{split}
 e^{\alpha u}z_k(u)
 &=(f*b^{*n})(u)\\
 &\leq C_\Phi(1+u)^a
       \int(1+v_-)^ab^{*n}(v)\,dv\\
 &\leq C_\Phi(1+u)^a n^aM_\alpha^{n-1}I_a^- .
 \end{split}
 \tag{PGD.17}\label{eq:pgd-tilt-proof}
\]
For the last inequality, in the product integral use
\((\sum_j t_j)_-\leq\sum_j(t_j)_-\) and
\[
 (1+\sum_{j=1}^n(t_j)_-)^a
 \leq\left(\sum_{j=1}^n[1+(t_j)_-]\right)^a
 \leq n^{a-1}\sum_{j=1}^n[1+(t_j)_-]^a.
\]
The last inequality is convexity of \(x^a\), with \(a>1\),
applied to the arithmetic mean of these \(n\) nonnegative numbers.
Integrating each of the \(n\) summands gives exactly
\(n^aM_\alpha^{n-1}I_a^-\), including \(n=1\).
Evenness of \(q_\Phi^2w\) and of \(w\) proves \(z_k(-u)=z_k(u)\),
so the positive half-line inequality proves (PGD.16) on the whole line.

The original Gamma density
\(\sigma(u)=|\Gamma(1/4+iu/2)|^2/(2\pi)\) has mass
\(\sqrt{2\pi}\) and satisfies the retained pointwise lower bound
\(\sigma(u)\geq c_\Gamma e^{-\alpha|u|}(1+|u|)^{-1/2}\).
It follows from (PGD.16) that
\[
 z_k(u)\leq F_k(1+|u|)^7\sigma(u),\qquad
 F_k=\frac{C_\Phi I_a^-}{c_\Gamma}
                    (k-1)^aM_\alpha^{k-2},\quad k\geq2.
 \tag{PGD.18}\label{eq:pgd-z-Gamma}
\]
The measure \(\sigma(u)du\) is used through this stated inequality;
it is not the \(k\)-fold Gamma reference of AT.4, and its mass
has not been raised to a different tensor power.

\subsection{All-vector bounds on the original growing windows}

Let \(\sigma\) retain that same density. Its complete monic
orthogonality calculation TG.7--TG.14 gives
\[
 \gamma_j=\sqrt{2\pi}(2j)!/4^j,\qquad
 u e_j=\sqrt{(j+1)(j+1/2)}e_{j+1}
           +\sqrt{j(j-1/2)}e_{j-1},
 \tag{PGD.19}\label{eq:pgd-Gamma-recurrence}
\]
where \(e_j\) are orthonormal, and the second term is absent at
\(j=0\). On degree at most \(n\), the two weighted shifts have
norms at most \(n+1\) and \(n\), as is seen by summing squared
coefficients in their orthogonal target frames. Hence
\(\|uf\|_\sigma\leq2(n+1)\|f\|_\sigma\).
For every integer \(r\geq0\), iterating through the degree increases
and expanding the nonnegative polynomial \((1+u^2)^r\) proves
\[
 \int(1+u^2)^r|f(u)|^2\sigma(u)\,du
       \leq[1+4(N+r)^2]^r\int|f(u)|^2\sigma(u)\,du,
 \quad\deg f\leq N.
 \tag{PGD.20}\label{eq:pgd-integer-multiplier}
\]
For \(r=0\) this is equality. For \(r\geq1\), the estimate
\(\|u^jf\|_\sigma\leq[2(N+r)]^j\|f\|_\sigma\),
\(0\leq j\leq r\), followed by the binomial theorem proves the
displayed constant. In particular, Holder's inequality with exponents
\(8/7\) and \(8\), followed by the case \(r=4\), gives
\[
 \begin{gathered}
 L_N=2^{7/2}[1+4(N+4)^2]^{7/2},\\
 \int(1+|u|)^7|f(u)|^2\sigma(u)\,du
                \leq L_N\int|f(u)|^2\sigma(u)\,du.
 \end{gathered}
 \tag{PGD.21}\label{eq:pgd-seven-multiplier}
\]
Explicitly \((1+|u|)^7\leq2^{7/2}(1+u^2)^{7/2}\), and
\[
 \int(1+u^2)^{7/2}|f|^2\sigma
 \leq\left(\int(1+u^2)^4|f|^2\sigma\right)^{7/8}
        \left(\int|f|^2\sigma\right)^{1/8}.
\]
Composition \(P\mapsto f(u)=P(k/2+iu)\) is the exact invertible
degree-preserving algebra map, with inverse
\(f\mapsto f((S-k/2)/i)\). Therefore these estimates apply to
every original polynomial vector, in the original sum coordinate.

Retain the original positive convolution lower envelope TW7:
\[
 \begin{gathered}
 B=42+8m,\quad M=B/2=21+4m,\qquad
 \vartheta_h=\int_{-1}^1w(t)\,dt>0,\\
 m_k(u)\geq c_h\vartheta_h^{k-3}
       e^{-\alpha(|u|+k-3)}(k-2+|u|)^{-B},\quad k\geq3,\\
 a_k=\frac{c_h\vartheta_h^{k-3}e^{-\alpha(k-3)}(k-2)^{-B}}
                   {C_\Gamma2^M},\qquad
 D_N=[1+4(N+M)^2]^M.
 \end{gathered}
 \tag{PGD.22}\label{eq:pgd-original-lower}
\]
Here \(c_h>0\) is exactly the source constant of the proved arithmetic
envelope AU.10, retained in CJ.15 and TW7. It depends on the fixed
full arithmetic packet and not on \(k,u,N\). The original proof of
that envelope, rather than a new lower bound on zeta, is the input in
(PGD.22). Its all-vector implication can be proved directly here.
The upper Gamma bound and
\((k-2+|u|)^B\leq(k-2)^B2^M(1+u^2)^M\) imply
\(m_k\geq a_k(1+u^2)^{-M}\sigma\).
Cauchy--Schwarz, followed by (PGD.20) at the integer \(r=M\), gives
\[
 \begin{split}
 \left(\int|f|^2\sigma\right)^2
 &\leq\left(\int(1+u^2)^{-M}|f|^2\sigma\right)
       \left(\int(1+u^2)^M|f|^2\sigma\right)\\
 &\leq D_N\left(\int(1+u^2)^{-M}|f|^2\sigma\right)
                   \left(\int|f|^2\sigma\right).
 \end{split}
\]
Division for \(f\ne0\), with the zero vector handled separately,
proves
\[
 \int|P(k/2+iu)|^2\sigma(u)\,du
 \leq\frac{D_N}{a_k}
           \int|P(k/2+iu)|^2m_k(u)\,du,\qquad \deg P\leq N.
 \tag{PGD.23}\label{eq:pgd-sigma-to-actual}
\]
Combining (PGD.8), (PGD.18), (PGD.21), and (PGD.23) now gives
the full finite-source comparison
\[
 \boxed{\quad
 c_kM_N\preceq W_N\preceq(c_k+E_{k,N})M_N,\qquad
 c_k=1+k^2A^2,\quad
 E_{k,N}=\frac{k^2F_kL_ND_N}{a_k},\quad k\geq3, N\geq0.
 \quad}
 \tag{PGD.24}\label{eq:pgd-full-window}
\]
It applies to arbitrary mixtures of all polynomial coefficients.
Restriction by multiplication by the literal \(\chi\) gives the same
constants on its full relation space. For \(N\geq q-1\), minimization
over each unchanged affine remainder fibre gives the same constants
for the two induced full \(q\)-dimensional quotient Grams. For smaller
\(N\), the same assertion concerns the image of the restricted
remainder map, with its actual dimension. None of these steps drops a relation source or its
phase energy. The original monomial matrices themselves are retained;
\(c_k\) is a proved lower comparison scalar, not a rescaling of them.

The logarithmic comparison constant and its full dependence are
\[
 \begin{gathered}
 \Lambda_{k,N}^{\rm ph}
     =\log\left(1+\frac{E_{k,N}}{c_k}\right),\qquad
 d_h=\alpha+\log(M_\alpha/\vartheta_h)>\alpha,\\
 \begin{aligned}
 \log\frac{E_{k,N}}{c_k}
 ={}&2\log k-\log(1+k^2A^2)+a\log(k-1)
       +(k-2)\log M_\alpha\\
    &-(k-3)\log\vartheta_h+\alpha(k-3)+B\log(k-2)\\
    &+\log\left(\frac{C_\Phi I_a^-C_\Gamma2^M}
                         {c_\Gamma c_h}\right)
       +\log L_N+\log D_N .
 \end{aligned}
 \end{gathered}
 \tag{PGD.25}\label{eq:pgd-log-cost}
\]
The strict inequality in \(d_h\) follows from
\(M_\alpha=\int\cosh(\alpha t)w(t)dt>\int_{-1}^1w(t)dt\):
the complement has positive measure with positive density. All terms
in the logarithm have been retained. For fixed \(h\),
\(\log L_N+\log D_N=O_h(1+\log(N+1))\), with polynomial
degree \(B+7=49+8m\) in \(N\). The other displayed terms prove,
uniformly for \(k\geq3,N\geq0\),
\[
 \Lambda_{k,N}^{\rm ph}
       =k d_h+O_h(\log k+\log(N+1)).
 \tag{PGD.26}\label{eq:pgd-cost-rate}
\]
To justify using \(\log(1+x)\) here, the displayed exact expression
for \(\log(E_{k,N}/c_k)\), together with \(A>0\) and the positive
lower bounds of \(L_N,D_N\), tends uniformly to \(+\infty\) as
\(k\to\infty\). Thus
\(\log(1+x)-\log x=\log(1+x^{-1})\) is uniformly bounded in
the asserted range; finitely many smaller \(k\) are absorbed by
the same fixed constant. This is a rate for the proved upper envelope,
not an assertion that the actual condition number has that growth.

\subsection{The exact signed four-endpoint consequence}

Retain \(\chi=\chi_{h,k}\),
\(q=[1+k(m-1)](k+1)^2\), and the original four endpoints
\(q-1,q,2q-1,2q\), with signs \(+,+,-,-\).
For any compatible positive source family \(X_N\), let
\[
 \begin{gathered}
 G_N(X)=(J_NX_N^{-1}J_N^*)^{-1},\\
 \mathcal B(X)=\log\det G_{q-1}(X)+\log\det G_q(X)
       -\log\det G_{2q-1}(X)-\log\det G_{2q}(X).
 \end{gathered}
 \tag{PGD.27}\label{eq:pgd-four-functional}
\]
The quotient is always by the same literal monic \(\chi\).
Its exact determinant formula is
\(\det G_N(X)=\det X_N/\det(B_N^*X_NB_N)\), including determinant
one for the zero-dimensional relation space. It follows by completing
the positive quadratic form on \(s_Nx+B_Ny\); the frame
\([s_N\ B_N]\) has determinant one by monic division.

For completeness the sharp dimension coefficient in the comparison
can be recovered on this very source. Interpolate the compatible
families by \(X_N(t)=(1-t)M_N+tW_N\). On the common space
\(\mathcal P_{2q}\), let \(P_N(t)\) be orthogonal projection onto
\(\mathcal P_N\), \(Q_N(t)\) projection onto
\(\chi\mathcal P_{N-q}\), and
\(C(t)=X_{2q}(t)^{-1}(W_{2q}-M_{2q})\).
Finite matrix differentiation of the preceding determinant identity
gives the exact, signed formula
\[
 \begin{gathered}
 \mathcal B(W)-\mathcal B(M)
       =\int_0^1\operatorname{Tr}((U(t)-V(t))C(t))\,dt,\\
 U=(Q_{2q-1}-Q_{q-1})+(Q_{2q}-Q_q),\qquad
 V=(P_{2q-1}-P_{q-1})+(P_{2q}-P_q).
 \end{gathered}
 \tag{PGD.28}\label{eq:pgd-exact-signed}
\]
All projections and adjoints here use the same actual \(X_{2q}(t)\)
metric. Each difference comes from nested subspaces. Reading their
orthogonal increments shows that both \(U\) and \(V\) have the
complete spectrum \(0^q,1^2,2^{q-1}\).
If \(c_1(t)\leq\cdots\leq c_{2q+1}(t)\) are the eigenvalues
of \(C(t)\), the trace against either positive operator with this
spectrum is between the sums obtained by pairing its increasing
eigenvalues with the decreasing or increasing lists \(c_j\).
One direct proof writes the operator as its range projection of rank
\(q+1\) plus its eigenvalue-two projection of rank \(q-1\).
In an eigenbasis of \(C\), a rank-\(r\) projection has diagonal
entries in \([0,1]\) summing to \(r\). Moving this fixed mass to
the \(r\) smallest or largest \(c_j\) proves the requisite two
trace bounds. Subtracting the complete sums therefore gives
\[
 |\operatorname{Tr}((U-V)C)|
 \leq c_{q+2}-c_q
       +2\sum_{j=1}^{q-1}(c_{2q+2-j}-c_j).
 \tag{PGD.29}\label{eq:pgd-spectral-bound}
\]
Let \(b_1\leq\cdots\leq b_{2q+1}\) be the positive generalized
eigenvalues of \(W_{2q}\) relative to \(M_{2q}\). Then
\(c_j(t)=(b_j-1)/(1-t+tb_j)\). This follows by diagonalizing
\(M_{2q}^{-1}W_{2q}\) in its original \(M_{2q}\) metric.
The fraction is increasing in \(b_j\) and integrates to
\(\log b_j\). Consequently (PGD.28)--(PGD.29) prove the exact
finite-spectrum and uniform bounds
\[
 \boxed{\quad
 |\mathcal B(W)-\mathcal B(M)|
 \leq\log\frac{b_{q+2}}{b_q}
      +2\sum_{j=1}^{q-1}\log\frac{b_{2q+2-j}}{b_j}
 \leq(2q-1)\Lambda_{k,2q}^{\rm ph}.
 \quad}
 \tag{PGD.30}\label{eq:pgd-signed-bound}
\]
The last inequality uses (PGD.24), so
\(c_k\leq b_j\leq c_k+E_{k,2q}\). Its total multiplicity is
\(1+2(q-1)=2q-1\). In particular the common scalar \(c_k\)
has cancelled in these ratios; none of the four original determinant
factors was changed before taking the ratio.

This source is itself a positive moment source. In fact
\(\mathcal B(W)>0\). To prove this directly, minimum sections at
two nested cutoffs satisfy
\(G_i(W)-G_j(W)=(R_i-R_j)^*W(R_i-R_j)\succeq0\), because
\(R_i-R_j\) is a relation and \(R_j\) is orthogonal to the
larger relation space. Thus each of the two logarithmic determinant
ratios in \(\mathcal B(W)\) is nonnegative. If their sum vanished,
both Gram differences, and hence both section differences, would
vanish. The ordered chain
\(G_{q-1}\succeq G_q\succeq G_{2q-1}\succeq G_{2q}\)
then forces all four Grams to agree, and the squared section-difference
identity forces all four sections to agree in the common source.
In particular the constant polynomial, which belongs to
\(\operatorname{im}R_{q-1}=\mathcal P_{q-1}\), would be orthogonal
to \(\chi\mathcal P_q\) at cutoff \(2q\). For
\(\chi(S)=\sum a_jS^j\), the polynomial
\(\chi^{\#_k}(S)=\sum\overline{a_j}(k-S)^j\) belongs to
\(\mathcal P_q\) and equals \(\overline{\chi(S)}\) on the
original line \(S=k/2+iu\). The asserted orthogonality would give
\(\int|\chi(k/2+iu)|^2r_k(u)du=0\), contrary to positivity of
the moment norm of this nonzero polynomial. This proof assigns no
sign to the difference \(\mathcal B(W)-\mathcal B(M)\).

Finally \(q=[1+k(m-1)](k+1)^2\) is the original unaltered degree,
so \(\log q=O_h(\log k)\). Equations (PGD.25)--(PGD.26) give
\[
 \frac{\Lambda_{k,2q}^{\rm ph}}{k}\longrightarrow d_h>\alpha,
 \qquad
 \frac{(2q-1)\Lambda_{k,2q}^{\rm ph}}{kq}
                                  \longrightarrow2d_h.
 \tag{PGD.31}\label{eq:pgd-amplified-rate}
\]
These are proved rates of this explicit estimate. They do not assert
that the actual signed correction is nonzero or attains the bound.
For the retained quartet spectral quantity
\(\mathcal L_k=2\delta[1+k(m-1)](k+1)
                 \lfloor(k+1)^2/4\rfloor\), one has
\(\mathcal L_k/(kq)\to\delta/2\). Thus the upper envelope in
(PGD.30), measured on that actual scale, tends to \(4d_h/\delta\).
Measured against \(4q\log k\), the same upper envelope diverges
as \(d_hk/(2\log k)\). This calculation preserves the exact
phase graph and all mixed terms while identifying the actual size of
the control furnished by its present decay estimate; it supplies
neither a shrinking exponent error nor a contradiction to the
counterfactual packet.
